\documentclass[11pt]{article}

\usepackage[margin=1in]{geometry}
\usepackage[T1]{fontenc}
\usepackage{amsmath,amssymb,bm}
\usepackage{booktabs}
\usepackage{array}
\usepackage{longtable}
\usepackage{enumitem}
\usepackage[hidelinks]{hyperref}

\newcommand{\ud}{\,\mathrm{d}}
\newcommand{\Grad}{\operatorname{Grad}}
\newcommand{\Dive}{\operatorname{Div}}
\newcommand{\dev}{\operatorname{dev}}
\newcommand{\tr}{\operatorname{tr}}
\newcommand{\diag}{\operatorname{diag}}
\newcommand{\bx}{\bm{x}}
\newcommand{\bX}{\bm{X}}
\newcommand{\bu}{\bm{u}}
\newcommand{\be}{\bm{e}}
\newcommand{\bN}{\bm{N}}
\newcommand{\bsig}{\bm{\sigma}}
\newcommand{\bP}{\bm{P}}
\newcommand{\bF}{\bm{F}}
\newcommand{\bB}{\bm{B}}
\newcommand{\bI}{\bm{I}}
\newcommand{\bR}{\bm{R}}
\newcommand{\bb}{\bm{b}}
\newcommand{\bw}{\bm{w}}
\newcommand{\Ione}{I_1}
\newcommand{\Itwo}{I_2}
\DeclareRobustCommand{\code}[1]{{\def\_{\textunderscore\allowbreak}\texttt{#1}}}
\DeclareRobustCommand{\file}[1]{{\def\_{\textunderscore\allowbreak}\texttt{#1}}}
\sloppy
\setlength{\emergencystretch}{2em}

\title{Continuum-Mechanics-MFEM\\[4pt]
\large Verification Manual (\texttt{apps/} and \texttt{tests/})}
\author{}
\date{September 2026}

\begin{document}
\maketitle

\begin{abstract}
\noindent
This manual documents how the flux-kernel framework of the root \file{src/}
library is verified.  It describes the application executable and the
scripts under \file{apps/}, the example meshes, and every input under
\file{apps/input/}: the problem it poses, the reference it is compared
with, the quantity checked and the tolerance, and the test that runs it.
The unit-level checks of the test suite (\file{tests/}) are summarised
first, since they verify the ingredients the application cases rely on.
The models and methods themselves are described in the theory manual
(\file{doc/theory\_manual.tex}); the closed-form solutions of the
homogeneous states are derived in \file{doc/incompressible\_hyperelasticity.tex}.
\end{abstract}

\tableofcontents

%============================================================================
\section{Organisation}
%============================================================================

\subsection{Levels of verification}

The checks form three tiers, each relying on the previous one.

\begin{enumerate}[nosep]
  \item \textbf{Unit and consistency tests} (\file{tests/test\_base},
    \file{test\_materials}): tensor algebra against hand values, dual
    derivatives against finite differences, the YAML schema and the
    expression parser, and every material's stress-free reference state,
    tangent, objectivity and small-strain limit.
  \item \textbf{Discretisation tests} (\file{tests/test\_solid\_mms},
    \file{test\_mixed}, \file{test\_homogeneous}, \file{test\_loading}):
    patch tests, assembled Jacobians against finite differences, manufactured
    solutions with measured convergence rates, homogeneous states against
    closed forms, and the boundary-condition and load-scheduling features.
  \item \textbf{Application cases} (\file{apps/input/}, driven by
    \file{tests/test\_verification}, \file{test\_benchmarks},
    \file{test\_loading}, \file{test\_parallel} and
    \file{apps/homogeneous\_compare.py}): complete inputs run through the
    same library wiring as the executable and compared with analytical
    solutions, literature benchmarks, frozen regression values, or a
    reference code.
\end{enumerate}

\subsection{Running the suite}

\begin{center}
\small
\begin{tabular}{@{}l p{11cm}@{}}
\toprule
Target & Content \\
\midrule
\code{make} & \file{build/libcmf.a}, \file{build/apps/solid\_mechanics}, \file{build/tests/*} \\
\code{make meshes} & regenerates \file{apps/mesh/*.msh} from the \file{.geo} sources with Gmsh (the \file{.msh} files are kept in the tree) \\
\code{make check} & serial, about 20 s: \file{test\_base}, \file{test\_materials}, \file{test\_solid\_mms}, \file{test\_mixed}, \file{test\_homogeneous}, \file{test\_loading} \\
\code{make homogeneous} & the executable on the homogeneous-deformation inputs, compared with the closed forms by \file{apps/homogeneous\_compare.py} \\
\code{make test} & everything: \code{check}, the Cook inputs serial and with four ranks, parallel consistency with two and four ranks, the incompressible Cook input, the homogeneous comparison, \file{test\_mixed --full}, \file{test\_benchmarks} serial and parallel, \file{test\_loading} with four ranks, \file{test\_verification}, and finally the Cook ratio gate \\
\bottomrule
\end{tabular}
\end{center}
Targets are run from the repository root, since inputs are referenced as
\file{apps/input/<set>/<case>.yaml}.  Every test returns a non-zero exit
status on the first failed check and prints the measured quantities next
to their references.

\subsection{The executable}

\file{apps/solid\_mechanics.cpp} contains no physics: it parses the input,
builds the parallel mesh, creates the problem selected by
\code{formulation}, installs the loads, and runs the quasi-static stepper
with a callback that writes ParaView cycles, probes and reactions after
every converged increment.  At startup it prints the mesh statistics and
the element and boundary attributes with their names; at the end the
Newton report, $\|\bu\|_{L^2}$, the internal energy, the probes and the
reactions.  The exit status is 0 on convergence, 2 if the load path did
not complete, and 1 on an input error.  Run as
\begingroup\small
\begin{verbatim}
./build/apps/solid_mechanics -i apps/input/finite_elasticity/cooks_membrane/cook.yaml
mpirun -np 4 ./build/apps/solid_mechanics -i <input.yaml>
\end{verbatim}
\endgroup

\subsection{Meshes}

Every input reads a Gmsh mesh from \file{apps/mesh/}, generated from a
\file{.geo} source with named physical groups that the boundary conditions
refer to.  The box-like geometries keep MFEM's Cartesian numbering (2D:
bottom 1, right 2, top 3, left 4; 3D: bottom $z{=}0$ 1, front $y{=}0$ 2,
right $x{=}L$ 3, back $y{=}L$ 4, left $x{=}0$ 5, top $z{=}L$ 6).

\begin{center}
\small
\begin{tabular}{@{}l p{3.3cm} p{6.6cm}@{}}
\toprule
Source & Meshes & Geometry and groups \\
\midrule
\file{square.geo} & \file{square.msh} & unit square, $n\times n$ quadrilaterals; bottom, right, top, left \\
\file{cook.geo} & \file{cook.msh} & Cook's membrane, transfinite quadrilaterals (the bilinear map of the former Cartesian input); bottom, right, top, left \\
\file{box.geo} & \file{cube}, \file{beam}, \file{cube10}, \file{shear\_cube}, \file{column\_*} & $L_x\times L_y\times L_z$ hexahedra; bottom, front, right, back, left, top \\
\file{cube\_tet.geo} & \file{cube\_tet.msh} & unit cube in tetrahedra, same groups \\
\file{annulus.geo} & \file{annulus.msh} & quarter annulus $A{=}1$, $B{=}2$, second order; bottom, outer, left, inner \\
\file{cylinder\_torsion.geo} & \file{cylinder\_torsion.msh} & solid cylinder $a{=}12.7$, $L{=}25.4$, second-order tetrahedra; bottom, top, wall \\
\file{plate\_hole.geo} & \file{plate\_hole.msh} & quarter plate $15\times10\times1$ with a 3\,mm hole, second order; xbot, ybot, xtop, ytop, hole, zbot, ztop \\
\file{plate\_hole\_2d.geo} & \file{plate\_hole\_2d.msh} & quarter of a $40\times40$ plate with a unit hole, second-order quadrilaterals \\
\file{tube\_quarter.geo} & \file{tube\_quarter.msh} & quarter tube $10/11\times5$, second order; xplane, yplane, zbot, ztop, inner, outer \\
\file{sphere\_octant.geo} & \file{sphere\_octant.msh} & octant shell $10/11$, second order; xplane, yplane, zplane, inner, outer \\
\file{footing.geo} & \file{footing.msh} & 50\,mm cube with a quarter patch on top; left, front, bottom, patch, top\_rest, right, back \\
\file{inclusion.geo} & \file{inclusion.msh} & octant of a cube with a spherical inclusion; volumes inclusion, matrix; the six box faces \\
\file{perturb\_column.py} & \file{column\_buckling}, \file{column\_euler} & adds the imperfection $x, y \mathrel{-}= \delta_0(1 - \cos 2\pi z/L)$ to a straight box mesh \\
\bottomrule
\end{tabular}
\end{center}

%============================================================================
\section{Unit and consistency tests}
%============================================================================

\subsection{\texttt{test\_base}}
Tensor algebra (determinant, inverse, contractions of rank 2 and 4) against
hand values in two and three dimensions; dual-number derivatives of scalar
composites and of $\det$, $\mathrm{inv}$ and tensor compositions against
Richardson-extrapolated central differences to $10^{-10}$; the YAML
schema: the documented example parses to the expected values, defaults for
omitted sections, malformed files give an error rather than a crash,
missing keys, wrong types, unknown keys and bad enumerations name the full
key path; boundary attributes by number and by physical-group name; load
schedules and their YAML forms (ramp, constant, table, components, step
segments, bisection control); the expression grammar (precedence,
right-associative power, functions, comparisons, parse errors with column
numbers); and material regions (inheritance, bulk-key replacement, errors).

\subsection{\texttt{test\_materials}}
For every model, at five random deformation gradients $\bF = \bI + 0.4\bR$
with $\det\bF > 0.3$:
\begin{itemize}[nosep]
  \item stress-free reference state $\bP(\bI) = \bm 0$ to machine precision;
  \item the dual-number tangent against central finite differences of
    $\bP$ (step $10^{-6}$), relative error $\le 10^{-6}$;
  \item objectivity $\bP(\bm Q\bF) = \bm Q\bP(\bF)$ for random rotations
    $\bm Q$ to $10^{-12}$;
  \item the small-strain limit: for $\bF = \bI + \epsilon\bm H$ the stress
    equals $\lambda\tr\bm\varepsilon\,\bI + 2\mu\bm\varepsilon$ with the
    resolved moduli ($\lambda = \kappa - \tfrac23\mu$ for the decoupled
    models);
  \item for the decoupled models, $\bP = \partial\Psi/\partial\bF$ and
    $\bP_{\mathrm{iso}} = \partial\Psi_{\mathrm{iso}}/\partial\bF$ by
    duals against finite differences, and the deviatoric property
    $\tr(\bP_{\mathrm{iso}}\bF^T) = 0$ at every $\bF$;
  \item Ogden reduces to the isochoric neo-Hookean model for
    $(\mu, \alpha) = (\mu, 2)$ and to Mooney--Rivlin for
    $(2c_1, 2), (-2c_2, -2)$; its analytic tangent matches finite
    differences also at coincident principal stretches;
  \item the plane-stress adapter: $P_{33} = 0$ and no out-of-plane shear
    stress, $J = 1$ for incompressible bases, the in-plane dual tangent
    against finite differences (the compressible branch differentiates a
    converged root), $\partial\Psi/\partial\bF = \bP$ in the plane
    (envelope theorem), objectivity under in-plane rotations, the
    small-strain plane-stress moduli $E/(1-\nu^2)$, $E\nu/(1-\nu^2)$ and
    the thickness strain $-\nu/(1-\nu)\,\varepsilon$, and uniaxial tension
    of an incompressible sheet, $P_{11} = \mu(\lambda - \lambda^{-2})$;
  \item the volumetric laws (quadratic, Simo--Taylor, logarithmic,
    $J\ln J$): $u(1) = u'(1) = 0$, $u''(1) = 1$, $u'$ and $u''$ against
    dual and finite-difference derivatives of $u$ and $u'$, and the
    isochoric neo-Hookean and Ogden models through the full decoupled
    checks above with each law; the factory applies the key and the
    coupled models refuse it;
  \item resolution of $(\mu, \kappa)$ from every admissible key combination
    and the refusal of the inadmissible ones.
\end{itemize}

%============================================================================
\section{Discretisation tests}
%============================================================================

\subsection{Patch tests, manufactured solutions, Jacobians (\texttt{test\_solid\_mms})}
On a perturbed quadrilateral mesh an affine displacement $\bu = \bm A\bX + \bm c$
prescribed on the boundary is reproduced exactly (patch test) for both
coupled models.  The smooth manufactured field
$\bu = \alpha(\sin\pi X\sin\pi Y,\ X^2Y(1-Y))$ with the body force
$\bb = -\rho_R^{-1}\Dive\bP(\bF(\bu))$ (divergence by central differences of
the closed-form stress) gives $L^2$ convergence rates of at least
$k + 0.9$ for $k = 1, 2$; on the finest $k = 2$ mesh the last Newton steps
above the round-off floor contract quadratically (order estimate
$\ge 1.5$); the assembled Jacobian matches central differences of the
residual along a random direction at a nonzero state.

\subsection{Mixed formulation (\texttt{test\_mixed})}
Affine patch tests with a constant pressure $p_0 = U'(\det(\bI + \bm A))$
for the decoupled materials at $\nu = 0.45$ and at $\kappa = \infty$, and
at $\nu = 0.45$ with the Simo--Taylor, logarithmic and $J\ln J$ volumetric
laws (exact to $10^{-14}$, quadratic Newton convergence); the block
Jacobian against finite differences ($4\times10^{-9}$), with the pressure
rows checked separately so that the law-dependent block $K_{pu}$ is
exercised; manufactured
solutions with a volume-preserving motion (the composition of two shears
$x = X + a\sin\pi Y$, $y = Y + b\sin\pi x$, $J = 1$, with the smooth
pressure $p = c\sin\pi X\cos\pi Y$) for the incompressible case, rates
$u \ge 3.0$, $p \ge 2.0$, and a compressible field at $\nu = 0.45$ with
$p = \kappa(J-1)$, rates $u \ge 3.4$, $p \ge 2.0$; and agreement of the
mixed and displacement formulations at finite $\kappa$ under refinement.
The pressure is made unique by a manufactured traction on one edge.
\code{--full} adds the $32\times32$ levels.

\subsection{Homogeneous deformations (\texttt{test\_homogeneous})}
For the six incompressible models with rubber-like parameters (isochoric
neo-Hookean $\mu = 1$; Mooney--Rivlin $c_1 = 0.4$, $c_2 = 0.1$; Yeoh
$0.5, -0.05, 0.01$; Gent $\mu = 1$, $J_m = 10$; Arruda--Boyce $\mu = 1$,
$N = 5$; Ogden $(0.63, 0.0012, -0.01)$, $(1.3, 5.0, -2.0)$):
\begin{enumerate}[nosep]
  \item at the material point, the code's Cauchy stress at
    $\bF = \diag(\lambda_a)$ with the analytic mean stress equals the
    analytic principal stresses of the plane-strain, uniaxial, equibiaxial
    and pure-shear states, and the deviatoric stress in simple shear equals
    its spectral closed form with $\sigma_{11} - \sigma_{22} = \gamma\sigma_{12}$;
  \item as finite element solutions of the mixed formulation on $4\times4$
    (plane strain) and $2\times2\times2$ (uniaxial) Q2/Q1 meshes, and of the
    plane-stress displacement formulation for the sheet states uniaxial,
    equibiaxial, pure shear and simple shear on a $4\times4$ Q2 mesh:
    displacement, pressure and every quadrature quantity (Cauchy and first
    Piola--Kirchhoff stress, $J$, von Mises, energy density, thickness
    stretch) in all presentations (both nodal projections, element
    averages, raw quadrature values), exact to solver tolerance with
    quadratic Newton convergence.
\end{enumerate}

\subsection{Boundary conditions and loading (\texttt{test\_loading})}
On Cook's membrane and unit boxes: staged loading (a traction ramped over
$[0, 0.5]$ and held, a second ramped over $[0.5, 1]$) equals the
single-stage solve, and the state at $t = 0.5$ equals the first traction
alone (path independence); a load-unload table returns to zero; with a
capped Newton the planned step fails and bisection recovers, while without
substepping the failure is reported; the tangent predictor reaches the
same converged state as the default start with fewer Newton iterations
(Cook's membrane with a prescribed edge, 18 to 9; a clamped sheared square
in the mixed formulation with the logarithmic volumetric law at
$\kappa/\mu = 1000$, 17 with a damped step to 6 full steps) and leaves a
traction-driven path unchanged; a programmatic time-dependent
coefficient with the constant schedule reproduces the ramp; the
manufactured problem written entirely as YAML expressions converges at
rate $\ge 2.8$ for $k = 2$, also with only the $x$ component prescribed on
two faces and the exact traction supplied for the other; a traction
expression linear in $t$ equals value plus ramp and one without $t$ is
assembled once; the uniaxial symmetry cube with rollers reproduces the
affine field and closed-form pressure of the full cube, with exactly the
expected number of essential degrees of freedom; rollers on all six faces
under a body force leave no rigid mode; the follower tangent matches
central differences on distorted quadrilateral and hexahedral meshes at a
random state; the follower pressure equals the dead pressure at $\bF = \bI$
and agrees to $O(p^2)$ under a small load; the thick-walled cylinder
inflates as Rivlin says (Section~\ref{sec:rivlincyl}); two-region
materials reproduce the single-material run when the parameters agree, have
a consistent tangent and deform less in the stiff half; and the reactions
equal $P_{11}$ times the face area on the uniaxial cube and balance the
applied resultant on Cook's membrane to $10^{-14}$.

\subsection{Parallel consistency (\texttt{test\_parallel})}
The Cook inputs of both formulations are run serially with tightened
tolerances (\code{--write}) and the recorded norms are matched to
$10^{-10}$ relative by runs with two and four ranks (\code{--check}).

%============================================================================
\section{Verification cases (\texttt{apps/input/finite\_elasticity/verification/})}
%============================================================================

Every input in this directory is compared with an independent reference
by a test.  The header of each input states the reference and the
formulas.  Unless noted, the material parameters are dimensionless.

\subsection{Homogeneous deformations of incompressible models}
\label{sec:homogeneous}

\paragraph{Problem.}
A unit square or cube carries an affine displacement $\bu = \bm G\bX$,
$\bm G = \diag(\lambda_a - 1)$, on its constrained faces, written as
expression strings (for example \code{["x", "-0.5*y"]}), with the remaining
faces free; the isochoric neo-Hookean model with \code{incompressible: true}
and $\mu = 1$ ($\mu = 0.5$\,MPa in the three uniaxial inputs, which share
their shear modulus with Section~\ref{sec:compuniaxial}).  The state is homogeneous, so
$\Dive\bP = 0$ holds identically and the affine field with a constant
pressure is the exact solution, which lies in the Taylor--Hood spaces.

\begin{center}
\small
\begin{tabular}{@{}l l l p{4.2cm}@{}}
\toprule
Input & Formulation & Stretches & Constrained faces \\
\midrule
\file{plane\_strain\_neo\_hookean} & mixed, 2D & $(2, \tfrac12, 1)$ & left, right \\
\file{plane\_stress\_neo\_hookean} & plane stress, 2D & $(2, 2^{-1/2}, 2^{-1/2})$ & left, right \\
\file{uniaxial\_neo\_hookean} & mixed, 3D & $(2, 2^{-1/2}, 2^{-1/2})$ & left, right \\
\file{symmetry\_uniaxial\_neo\_hookean} & mixed, 3D & same, octant model & rollers on $X, Y, Z = 0$; $u_x$ on right \\
\file{incompressible\_uniaxial\_iso\_neo\_hookean} & mixed, 3D & $(1.5, 1.5^{-1/2}, 1.5^{-1/2})$, octant model & rollers on $X, Y, Z = 0$; $u_x$ on right \\
\file{equibiaxial\_neo\_hookean} & mixed, 3D & $(1.5, 1.5, 1.5^{-2})$ & left, right, front, back \\
\file{pure\_shear\_neo\_hookean} & mixed, 3D & $(1.5, 1, 1.5^{-1})$ & left, right, front, back \\
\bottomrule
\end{tabular}
\end{center}

\paragraph{Reference.}
The closed forms of \file{incompressible\_hyperelasticity.tex}, Sections
4.1--4.3 and 5.1--5.3.  With $\beta_a = \lambda_a\,\partial\Psi/\partial\lambda_a$
up to an isotropic term ($\beta_a = 2\Psi_1\lambda_a^2 - 2\Psi_2\lambda_a^{-2}$
for invariant models, $\sum_r\mu_r\lambda_a^{\alpha_r}$ for Ogden), the
principal Cauchy stresses are $\sigma_a = \beta_a - \beta_f$ where $f$ is
the traction-free axis, $p = \tfrac13\sum\sigma_a$,
$\sigma_{vm} = \sqrt{\tfrac32\sum(\sigma_a - p)^2}$, $P_a = \sigma_a/\lambda_a$.
For the neo-Hookean model: uniaxial $\sigma_1 = \mu(\lambda^2 - \lambda^{-1})$;
equibiaxial $\sigma_1 = \sigma_2 = \mu(\lambda^2 - \lambda^{-4})$; pure shear
$\sigma_1 = \mu(\lambda^2 - \lambda^{-2})$, $\sigma_2 = \mu(1 - \lambda^{-2})$;
plane-strain extension $\sigma_{11} = \mu(\lambda^2 - \lambda^{-2})$,
$\sigma_{33} = \mu(1 - \lambda^{-2})$.

\paragraph{Check.}
\file{apps/homogeneous\_compare.py} (\code{make homogeneous}) runs every
input, reads the affine data and the constrained faces from the YAML to
derive the stretches and the free axis, and compares the probed
displacement, pressure or thickness stretch, von Mises stress, $J$ and
every component of $\bsig$ and $\bP$, in the nodal and the element
presentations, with the closed forms; relative tolerance $10^{-8}$
(measured agreement about $10^{-13}$).  Any of the six models runs in the
same inputs by changing the \code{material} line; the script's closed
forms cover all of them.

\subsection{Compressible uniaxial tension}
\label{sec:compuniaxial}

\paragraph{Problem.}
The unit cube with rollers on $X, Y, Z = 0$ and the face $X = 1$ moved by
$0.5\,t$ in $x$ only (stretch $\lambda_1 = 1.5$), all other faces free.
Six inputs with the same small-strain moduli, $\mu = 0.5$\,MPa and
$\nu = 0.45$ ($\kappa = 29/6$\,MPa, $\kappa/\mu \approx 9.7$, compressible
enough for the volumetric laws to differ): coupled neo-Hookean and
Saint Venant--Kirchhoff with $E = 1.45$, $\nu = 0.45$ (displacement
formulation), and the isochoric neo-Hookean model with $\mu = 0.5$,
$\kappa = 29/6$ in the mixed formulation with each volumetric law:
quadratic (\file{\ldots\_mixed}), \code{simo\_taylor}, \code{j\_log\_j} and
\code{logarithmic} (the mixed constraint with $u'' \ne 1$).
\file{incompressible\_uniaxial\_iso\_neo\_hookean}
(Section~\ref{sec:homogeneous}) is their limit $\kappa = \infty$.

\paragraph{Reference.}
The state is homogeneous with $\bF = \diag(\lambda_1, \lambda_2, \lambda_2)$
and the lateral stretch fixed by $\sigma_{22} = 0$.  The test solves
$P_{22}(\diag(\lambda_1, \lambda_2, \lambda_2)) = 0$ for $\lambda_2$ with
the material's own stress function (which carries the volumetric law), then
forms $P_{11}$, $\sigma_{11}$ and $J$.

\paragraph{Check} (\file{test\_verification}).  Axial displacement to
$10^{-10}$, lateral displacements to $10^{-8}$, $P_{11}$, $\sigma_{11}$ and
$J$ to $10^{-7}$ relative, and the lateral nominal stresses vanish to
$10^{-7}P_{11}$.

\paragraph{Large stretch} (\file{apps/uniaxial\_plots.py},
\file{apps/neo\_hookean\_compare.py}).  The first script drives every
uniaxial input of this and the previous section to $\lambda_1 = 8$ in 70
load steps (Saint Venant--Kirchhoff to 0.95 of $\sqrt{1 + 1/\nu}$, where its
lateral stretch $\lambda_2^2 = 1 - \nu(\lambda_1^2 - 1)$ vanishes), reads the
ParaView series back and plots $P_{11}(\lambda_1)$ on an analytical curve
coded independently of \file{src/}: $\lambda_2$ by bisection of
$P_{22} = 0$ with
$P_{aa} = \mu(\lambda_a - \lambda_a^{-1}) + \lambda\ln J/\lambda_a$ (coupled
neo-Hookean) or
$P_{aa} = \mu J^{-2/3}(\lambda_a - I_1/(3\lambda_a)) + U'(J)J/\lambda_a$
(decoupled), the closed forms $\lambda_2^2 = 1 - \nu(\lambda_1^2 - 1)$,
$P_{11} = E\lambda_1(\lambda_1^2 - 1)/2$ for Saint Venant--Kirchhoff, and
$P_{11} = \mu(\lambda_1 - \lambda_1^{-2})$ when incompressible.  Measured
agreement is $3\times10^{-12}$ relative or better at every step, and the
nodal $P_{11}$ is uniform to $2\times10^{-11}$.  The second script overlays
the neo-Hookean responses: at $\lambda_1 = 8$ the decoupled models lie 33\%
(quadratic) to 51\% (logarithmic, $J = 2.9$, beyond the peak of $\ln J/J$
at $J = e$) below the incompressible curve, while the coupled model, for
which $P_{22} = 0$ gives $P_{11} = \mu(\lambda_1 - \lambda_2^2/\lambda_1)$,
stays on it.

\subsection{Manufactured solutions}

\paragraph{Problem.}
An exact displacement is prescribed on the whole boundary and the body
force $\bb = -\rho_R^{-1}\Dive\bP(\bF(\bu))$ that makes it the solution is
supplied as expression strings generated by sympy (expanded, no
simplification).  Saint Venant--Kirchhoff with $E = 250$, $\nu = 0.3$,
$\rho_R = 2.5$ unless noted, $k = 2$.

\begin{center}
\small
\begin{tabular}{@{}l l p{7.8cm}@{}}
\toprule
Input & Mesh & Exact displacement \\
\midrule
\file{mms\_2d\_plane\_strain} & \file{square.msh} & $0.05\,(\sin\pi x\sin\pi y,\ x^2y(1-y))$ \\
\file{mms\_3d\_hex} & \file{cube.msh} & $0.05\,(\sin\pi x\sin\pi y\sin\pi z,\ x^2y(1-y)z,\ xz(1-z)\cos\tfrac{\pi y}{2})$ \\
\file{mms\_3d\_tet} & \file{cube\_tet.msh} & the same field on tetrahedra \\
\file{mms\_3d\_mixed} & \file{cube.msh} & $0.05\,(\sin\pi y\sin\pi z,\ z(1-z),\ 0)$; isochoric neo-Hookean, $\mu = 40$, $\kappa = 400$, mixed formulation \\
\bottomrule
\end{tabular}
\end{center}
The mixed field has a strictly upper-triangular gradient, so $J = 1$
exactly and the exact pressure $\kappa(J - 1)$ vanishes; it is applied in
four increments because on fine meshes the exact boundary data applied at
once inverts the elements next to the boundary while the interior is still
at rest.

\paragraph{Check.}
$L^2$ displacement error under \code{serial\_refine} 1, 2, 3 (two
dimensions, \file{test\_loading}) or 1, 2 (three dimensions,
\file{test\_verification}); the rate $\log_2(e_{h}/e_{h/2})$ must be at
least 2.8, i.e.\ third order for $k = 2$.

\subsection{Rivlin's inflation of a thick-walled cylinder}
\label{sec:rivlincyl}

\paragraph{Problem.}
Quarter annulus $A = 1$, $B = 2$ (\file{annulus.msh}, second-order
geometry), plane strain, incompressible neo-Hookean with $\mu = 1$, mixed
formulation; rollers on the two symmetry lines through component-wise
Dirichlet data; follower pressure $P = 0.3$ per unit current area on the
inner surface, ramped over the load path.

\paragraph{Reference} (Rivlin; Ogden, \emph{Non-linear Elastic
Deformations}, Section 5.3.1).  With $a, b$ the deformed radii and
$b^2 = B^2 - A^2 + a^2$,
\begin{equation}
  P = \mu\Bigl[\ln\frac{B}{A} - \ln\frac{b}{a} + \frac{A^2 - a^2}{2}\Bigl(\frac1{b^2} - \frac1{a^2}\Bigr)\Bigr];
  \qquad P = 0.3 \ \Rightarrow\ a = 1.28158,\ b = 2.15463 .
\end{equation}

\paragraph{Check} (\file{test\_loading}).  On the base mesh and one
refinement: inner and outer radii within $0.2\%$ of the inflation
$(a - A)$, $(b - B)$; the radial displacement is the same at $(A, 0)$ and
$(0, A)$ to $10^{-6}$; the rollers hold the symmetry lines to $10^{-12}$;
the final Newton solve takes at most six iterations (quadratic convergence
with the follower tangent); and the error decreases under refinement.  The
README quotes the agreement as $10^{-6}$ on the refined mesh.

\subsection{Green--Zerna inflation of a thick-walled sphere}

\paragraph{Problem.}
Octant of a spherical shell $A = 10$, $B = 11$ (\file{sphere\_octant.msh},
curved second-order tetrahedra), incompressible neo-Hookean $\mu = 1$,
mixed formulation, rollers on the three symmetry planes, follower pressure
$P = 0.100975807337556$ on the inner surface.

\paragraph{Reference} (Green and Zerna; Ogden).  With $\lambda_a = a/A$,
$\lambda_b = b/B$ and $b^3 = B^3 - A^3 + a^3$,
\begin{equation}
  P = \frac{\mu}{2}\Bigl[\Bigl(\frac4{\lambda_b} + \frac1{\lambda_b^4}\Bigr) - \Bigl(\frac4{\lambda_a} + \frac1{\lambda_a^4}\Bigr)\Bigr],
\end{equation}
and the chosen $P$ gives $\lambda_a = 1.2$ exactly ($a = 12$,
$b = 12.72190$), below the pressure maximum of this shell near
$\lambda_a = 1.4$.  The test inverts the relation by bisection.

\paragraph{Check} (\file{test\_verification}).  Inner and outer radii
within $0.5\%$ of the inflation; the inner radius on the $z$ axis equals
that on the $x$ axis to $10^{-3}(a - A)$ (spherical symmetry).

\subsection{Rivlin's torsion of a solid cylinder}

\paragraph{Problem.}
Solid cylinder $a = 12.7$, $L = 25.4$ (\file{cylinder\_torsion.msh}, curved
second-order tetrahedra), incompressible neo-Hookean $\mu = 1$, mixed
formulation.  The bottom face is clamped and the top face rotated rigidly
by $\theta = 1$\,rad with $u_z = 0$ ($\bu = \bR(\psi z)\bX - \bX$ on the
face, twist $\psi = \theta/L$, rim shear $\psi a = 0.5$) in 20 increments; a
larger first increment leaves the elements under the rotated face inverted
before Newton starts.

\paragraph{Reference.}
Rivlin's universal solution holds exactly: $\bu = \bR(\psi z)\bX - \bX$
throughout, the lateral surface is traction free by itself, and in
cylindrical components with $\gamma = \psi r$
\begin{gather}
  \sigma_{rz} = 0, \quad \sigma_{\theta z} = \mu\psi r, \quad
  \sigma_{rr} = \sigma_{zz} = -\tfrac12\mu\psi^2(a^2 - r^2), \quad
  \sigma_{\theta\theta} = \sigma_{rr} + \mu\psi^2r^2, \\
  p = \sigma_{rr} + \tfrac13\mu\psi^2r^2, \qquad
  \sigma_{vm} = \mu\sqrt{\gamma^4 + 3\gamma^2}.
\end{gather}

\paragraph{Check} (\file{test\_verification}).  Displacement field to
$2\times10^{-3}$; on the mid-plane the shear stress $\sigma_{\theta z}$
and the von Mises stress to $2\%$, the normal stresses and the pressure to
$3\%$ of the stress scale, and the absence of $r\theta$ and $rz$ shear to
$2\%$ (curved P2 tetrahedra).

\subsection{Euler--Bernoulli cantilever in the linear limit}

\paragraph{Problem.}
A $10\times1\times1$ beam (\file{beam.msh}, $20\times2\times2$ hexahedra,
$k = 2$) clamped at $x = 0$ with a downward shear traction on $x = L$ small
enough that $\max|\Grad\bu| \approx 5\times10^{-5}$; compressible
neo-Hookean, and the same problem with Saint Venant--Kirchhoff.  The
systems AMG options are selected in the input because the elasticity
options stall on this slender mesh; the test also checks that the
automatic fallback converges.

\paragraph{Reference.}
Euler--Bernoulli tip deflection $PL^3/(3EI) = 3.2\times10^{-4}$.

\paragraph{Check} (\file{test\_benchmarks}).  $\max|\Grad\bu| \le 10^{-4}$
(the load is in the linear regime); tip deflection within $0.15\%$ of the
beam value; the two materials agree in the linear limit.

\subsection{Euler buckling of a clamped column (Southwell plot)}

\paragraph{Problem.}
A $1\times1\times20$ column (\file{column\_euler.msh}, $2\times2\times40$
hexahedra) carrying the first-mode imperfection
$x, y \mathrel{-}= 0.0025\,(1 - \cos 2\pi z/20)$; compressible neo-Hookean
$E = 100$, $\nu = 0.3$.  The base is clamped, the top held in $x, y$ and
shortened by $u_z = -0.12\,t$, an axial strain up to $0.6\%$, about $73\%$
of the Euler strain of the clamped-clamped column,
$\varepsilon_{cr} = 4\pi^2I/(AL^2) = 8.22\times10^{-3}$.

\paragraph{Reference.}
Below the limit point the mid-height lateral deflection of an imperfect
column grows as $u = \delta_0\,(P/P_{cr})/(1 - P/P_{cr})$ (Southwell), so
$1/u$ is linear in $1/\varepsilon$ with slope $\varepsilon_{cr}/\delta_0$
and intercept $-1/\delta_0$.

\paragraph{Check} (\file{test\_verification}).  A least-squares fit of
$1/u$ against $1/\varepsilon$ over the load steps recovers
$\varepsilon_{cr}$ within $5\%$ of the Euler value and the imperfection
amplitude $\delta_0$ within $15\%$ of the geometric one (0.005).

\subsection{Kirsch's plate with a circular hole}

\paragraph{Problem.}
Quarter of a $40\times40$ plate with a hole of radius $a = 1$
(\file{plate\_hole\_2d.msh}, quadrilaterals of size 0.04 at the hole,
curved second-order geometry), plane stress, compressible neo-Hookean
$E = 1000$, $\nu = 0.3$; rollers on the symmetry lines, far-field traction
$\sigma_0 = 0.1$ on $x = 20$ (strain $10^{-4}$, linear regime).

\paragraph{Reference.}
Kirsch's infinite-plate solution: $\sigma_{xx} = 3\sigma_0$ at the hole top
$(0, a)$, $\sigma_{yy} = -\sigma_0$ at the hole side $(a, 0)$, and along
the $y$ axis $\sigma_{xx}/\sigma_0 = 1 + a^2/(2y^2) + 3a^4/(2y^4)$.  The
finite width $W = 20a$ raises the concentration by a few tenths of a
percent (Howland).

\paragraph{Check} (\file{test\_verification}).  Probed Cauchy stresses
$\sigma_{xx}$, $\sigma_{yy}$ within $2\%$ of $3\sigma_0$ at the probe
points, and no shear stress on the symmetry axes to $1\%$ of $3\sigma_0$.

%============================================================================
\section{Benchmarks (\texttt{apps/input/finite\_elasticity/cooks\_membrane/})}
%============================================================================

\paragraph{Problem.}
Cook's membrane: the quadrilateral with corners $(0,0)$, $(48,44)$,
$(48,60)$, $(0,44)$ (\file{cook.msh}), plane strain, left edge clamped,
uniform upward shear traction on the right edge, $k = 2$.  Three inputs:
\begin{center}
\small
\begin{tabular}{@{}l l p{5.6cm} l@{}}
\toprule
Input & Formulation & Material & Traction (resultant) \\
\midrule
\file{cook} & displacement & neo-Hookean $E = 250$, $\nu = 0.3$ & 3.75 (60) \\
\file{cook\_incompressible} & mixed Q2--Q1 & isochoric neo-Hookean $\mu = 80.194$, $\kappa = \infty$ & 6.25 (100) \\
\file{cook\_nearly\_incompressible} & mixed Q2--Q1 & the same with $\nu = 0.4999$ ($\kappa = 4\times10^5$) & 6.25 (100) \\
\bottomrule
\end{tabular}
\end{center}

\paragraph{Reference.}
There is no closed form; the benchmark is a literature standard for
locking and mesh convergence.  The test performs a self-convergence study
(the base $4\times4$ mesh and four uniform refinements up to $64\times64$
for the compressible case, up to $32\times32$ for the incompressible one)
and holds the finest corner displacement as a frozen regression value.

\paragraph{Check} (\file{test\_benchmarks}).  Monotone self-convergence of
the vertical displacement of the top-right corner $(48, 60)$; the frozen
values $u_y = 4.905891700497$ (compressible, $64\times64$; $30.7\%$ of the
16\,mm edge) and $u_y = 6.930412595013$ (incompressible, $32\times32$;
$43.3\%$) reproduced to $10^{-8}$ relative, serially and with four ranks.
The final step of \code{make test}, \code{--cook-ratio-gate}, asserts the
plan's requirement that successive differences shrink by at least three
per refinement.  That threshold is \emph{not} met (measured ratios 2.34,
2.45, 2.31): uniform refinement is limited by the singularity at the
108-degree clamped-free corner, so the point value converges at about
$h^{1.3}$.  Everything before that step is green; the threshold is kept
as written rather than relaxed, so \code{make test} ends with this known
failure.

%============================================================================
\section{Validation against a reference code: Anand\textquotesingle s examples}
\label{sec:anand}
%============================================================================

The inputs of \file{apps/input/anand\_coupled\_theories/finite\_elasticity/}
reproduce a published set of finite element examples.

\paragraph{Reference.}
The ten "Finite Elasticity" examples of Lallit Anand's \emph{Introduction
to Coupled Theories in Solid Mechanics} (Oxford University Press, 2025),
from the FEniCSx companion codes by Eric Stewart and Lallit Anand
(\url{solidmechanicscoupledtheories.github.io}).  All use Arruda--Boyce
with $G_0 = 280$\,kPa, $\lambda_L = 5.12$ ($N = \lambda_L^2 = 26.2144$) and
$K = 1000\,G_0$, in kPa and mm, the reference's logarithmic volumetric law
$p = K\ln J/J$ (\code{volumetric: logarithmic}, run with
\code{solver.predictor: tangent}), the mixed formulation on
Q2--Q1 or P2--P1 elements, and the reference's geometry, boundary
conditions, load histories and step counts.  Every input probes the points of the
reference's plots after every step and prints the reactions, so the curves
can be read from the log.  These are validation cases: the comparison is
with another code's results and with analytical curves where they exist,
not a test with a tolerance.

\begin{center}
\small
\begin{tabular}{@{}l l p{8.6cm}@{}}
\toprule
Input & Reference & Problem \\
\midrule
\file{01\_uniaxial\_tension} & 3D01 & 10\,mm cube, rollers on three planes, stretch 7.75 in $y$ over 100 steps \\
\file{02\_simple\_shear} & 3D02 & 1\,mm cube, base clamped, top face $u_x = \sin 4\pi t$ (two cycles of shear strain 1) \\
\file{03\_cylinder\_torsion} & 3D03 & cylinder $R = 12.7$, $L = 25.4$, top face rotated rigidly by 2.5\,rad with $u_z = 0$, 20 steps \\
\file{04\_plate\_with\_hole} & 3D04 & quarter plate $15\times10\times1$ with a 3\,mm hole, stretch 3 in $x$ \\
\file{05\_cylinder\_inflation} & 3D05 & quarter tube $10/11\times5$, plane-strain rollers, follower pressure to 50\,kPa \\
\file{06\_sphere\_inflation} & 3D06 & octant shell $10/11$, follower pressure to 35\,kPa \\
\file{07\_cube\_footing} & 3D07 & 50\,mm cube, follower pressure 1500\,kPa on a quarter patch of the top, held in $x, y$ there \\
\file{08\_column\_buckling} & 3D08 & $1\times1\times20$ column with the first-mode imperfection, shortened by 2.5\,mm \\
\file{09\_spherical\_inclusion} & 3D09 & octant of a cube with a ten times stiffer spherical inclusion (\code{material.regions}), stretch 2 \\
\file{10\_column\_twist} & 3D10 & $1\times1\times3$ column, top face turned through $2\pi$ \\
\bottomrule
\end{tabular}
\end{center}

\paragraph{Differences from the reference that change the numbers.}
The Arruda--Boyce model here is the five-term series in $\Ione/N$ rather
than the Pad\'e approximation of the inverse Langevin function, so it is
softer near the locking stretch (visible in case 01 above a stretch of
about 3); case 02 keeps the quadratic law $p = K(J-1)$, because with the
logarithmic law the run fails at $t = 0.064$ for any increment size: the
tangent bulk modulus $K(1 - \ln J)/J^2$ of that law softens in dilatation
($0.18\,K$ at $J = 1.67$, zero at $J = e$), and at the clamped top corners,
stress singularities that this mesh and quadrature resolve sharply, $J$
grows with the shear until the volumetric response is lost; hexahedra
replace tetrahedra on the boxes; and a failed
increment is bisected instead of ending the run (cases 05 and 06 stop early
in the reference).

\paragraph{Plots.}
\file{apps/anand\_plots.py} reproduces the result plots of the reference
pages from the logs (save the stdout of each run as
\file{out/anand\_coupled\_theories/finite\_elasticity/logs/<case>.log}, or
pass \code{--logs}) and overlays an analytical curve where one exists:
the homogeneous incompressible and nearly incompressible Arruda--Boyce
response for the uniaxial and shear blocks (with both this code's series
and the reference's Pad\'e form), Rivlin's universal torsion for torque and
axial force by quadrature in $r$, the incompressible thick-walled cylinder
and sphere inflation by quadrature with their limit points, the Euler load
$4\pi^2EI/L^2$ for the clamped column, and the matrix-only curve for the
inclusion.

\paragraph{What the plots show.}
The uniaxial cube, the torsion and both inflations lie on their reference
curves (the uniaxial one on the nearly incompressible solution at
$K = 1000\,G$, with the incompressible limit a few percent above at the
largest stretch).  The sheared block carries about $17\%$ less nominal
shear stress than homogeneous simple shear at a shear strain of 1, because
its lateral faces are free where simple shear needs tractions, and its two
cycles retrace one curve (elastic, no hysteresis).  The sphere stops at its
limit pressure of 34\,kPa and the cylinder near 38\,kPa, where the
reference also stopped.  The buckling column reaches 7.3\,mN at 0.2\,mm of
shortening, $5\%$ above the Euler load of the clamped column, and then
rises slowly to 7.5\,mN at 2.5\,mm along the hardening post-buckling path
of the elastica (the Euler value uses $E = 3G$ and neglects the finite
section).  The twisted column needs a compressive axial force of 78\,mN to
keep its length over a full turn (Poynting effect), as does the torsion
cylinder.

%============================================================================
\section{Summary}
%============================================================================

\begin{center}
\small
\renewcommand{\arraystretch}{1.15}
\begin{tabular}{@{}>{\raggedright\arraybackslash}p{3.0cm} >{\raggedright\arraybackslash}p{3.9cm} >{\raggedright\arraybackslash}p{4.4cm} >{\raggedright\arraybackslash}p{3.6cm}@{}}
\toprule
Case & Reference & Check & Test \\
\midrule
homogeneous deformations & closed forms (six models) & all probed fields, $10^{-8}$ & \file{homogeneous\_compare.py} \\
compressible uniaxial & lateral stretch from $P_{22} = 0$ & $u$, $P_{11}$, $\sigma_{11}$, $J$ to $10^{-7}$ & \file{test\_verification} \\
manufactured solutions & exact field, sympy body force & $L^2$ rate $\ge 2.8$ ($k = 2$) & \file{test\_loading}, \file{test\_verification} \\
cylinder inflation & Rivlin & radii to $0.2\%$, refinement & \file{test\_loading} \\
sphere inflation & Green--Zerna & radii to $0.5\%$, symmetry & \file{test\_verification} \\
torsion & Rivlin universal solution & $u$ to $2\times10^{-3}$, stresses 2--3\% & \file{test\_verification} \\
cantilever & Euler--Bernoulli & tip deflection to $0.15\%$ & \file{test\_benchmarks} \\
column buckling & Euler load, Southwell & $\varepsilon_{cr}$ to $5\%$ & \file{test\_verification} \\
plate with hole & Kirsch & stresses to $2\%$ & \file{test\_verification} \\
Cook's membrane & frozen values, self-convergence & $10^{-8}$; ratio gate not met & \file{test\_benchmarks} \\
Anand examples & FEniCSx codes, analytical curves & plots & \file{anand\_plots.py} \\
parallel consistency & serial run & norms to $10^{-10}$ & \file{test\_parallel} \\
\bottomrule
\end{tabular}
\end{center}

%============================================================================
\small\begin{thebibliography}{9}
%============================================================================
\bibitem{ogden} R.\,W. Ogden, \emph{Non-Linear Elastic Deformations}, Ellis Horwood, 1984 (Dover, 1997).
\bibitem{rivlin} R.\,S. Rivlin, Large elastic deformations of isotropic materials IV, \emph{Phil. Trans. R. Soc. A} 241 (1948) 379--397.
\bibitem{greenzerna} A.\,E. Green, W. Zerna, \emph{Theoretical Elasticity}, 2nd ed., Oxford University Press, 1968.
\bibitem{kirsch} E.\,G. Kirsch, Die Theorie der Elastizit\"at und die Bed\"urfnisse der Festigkeitslehre, \emph{Z. VDI} 42 (1898) 797--807.
\bibitem{howland} R.\,C.\,J. Howland, On the stresses in the neighbourhood of a circular hole in a strip under tension, \emph{Phil. Trans. R. Soc. A} 229 (1930) 49--86.
\bibitem{southwell} R.\,V. Southwell, On the analysis of experimental observations in problems of elastic stability, \emph{Proc. R. Soc. A} 135 (1932) 601--616.
\bibitem{cook} R.\,D. Cook, Improved two-dimensional finite element, \emph{J. Struct. Div. ASCE} 100 (1974) 1851--1863.
\bibitem{anand} L. Anand, \emph{Introduction to Coupled Theories in Solid Mechanics}, Oxford University Press, 2025; companion codes by E. Stewart and L. Anand, \url{solidmechanicscoupledtheories.github.io}.
\bibitem{ab} E.\,M. Arruda, M.\,C. Boyce, A three-dimensional constitutive model for the large stretch behavior of rubber elastic materials, \emph{J. Mech. Phys. Solids} 41 (1993) 389--412.
\end{thebibliography}

\end{document}
